% ============================================================
%  Análisis de imagen: Poder y Protesta
%  Escultura de Banksy, Waterloo Place, Londres (2026)
% ============================================================
\documentclass[12pt, a4paper]{article}

% --- Paquetes ---
\usepackage[utf8]{inputenc}
\usepackage[T1]{fontenc}
\usepackage[spanish]{babel}
\usepackage{geometry}
\geometry{top=3cm, bottom=3cm, left=3cm, right=3cm}

\usepackage{setspace}
\onehalfspacing

\usepackage{parskip}
\setlength{\parindent}{0.5cm}

\usepackage{csquotes}
\usepackage{microtype}
\usepackage{hyperref}
\hypersetup{
    colorlinks=true,
    linkcolor=black,
    urlcolor=blue,
    citecolor=black
}

% Encabezado/pie de página
\usepackage{fancyhdr}
\pagestyle{fancy}
\fancyhf{}
\rhead{\small Poder y Protesta}
\lhead{\small Análisis de imagen}
\cfoot{\thepage}

% --- Título ---
\title{
    \vspace{-1cm}
    \large \textbf{Análisis de imagen: Poder y Protesta} \\[0.4cm]
    \normalsize \textit{La escultura de Banksy en Waterloo Place, Londres (2026)} \\[0.3cm]
    \small A partir del capítulo ``Poder y Protesta'' de Peter Burke \\
    en \textit{Visto y no visto. El uso de la imagen como documento histórico}
}
\author{}
\date{}

% ============================================================
\begin{document}

\maketitle
\thispagestyle{fancy}
\vspace{-1.5cm}

% ---- Ficha de la imagen ----
\section*{Ficha de la imagen}

\begin{tabular}{p{3.5cm} p{10.5cm}}
\textbf{Título (oficioso):} & \textit{Blind Patriotism} [Patriotismo ciego] \\[4pt]
\textbf{Autor:}             & Banksy (identidad no confirmada oficialmente; Reuters, en marzo de 2026, señaló al artista Robin Gunningham, sin que Banksy lo confirme ni lo niegue) \\[4pt]
\textbf{Fecha:}             & Instalada en la madrugada del 29 de abril de 2026; confirmada por el artista el 30 de abril de 2026 mediante un vídeo publicado en su cuenta de Instagram \\[4pt]
\textbf{Lugar:}             & Waterloo Place, Westminster, Londres (Reino Unido) \\[4pt]
\textbf{Técnica/material:} & Escultura de resina, sobre plinto de concreto; firma ``Banksy'' grabada en la base \\[4pt]
\textbf{Fuentes:}          & NPR (2 de mayo de 2026): \url{https://www.npr.org/2026/05/02/nx-s1-5808871/banksy-statue-london}; CNN (1.º de mayo de 2026): \url{https://www.cnn.com/2026/05/01/style/banksy-new-statue-london-intl-hnk}; Wikipedia: \url{https://en.wikipedia.org/wiki/Banksy_statue}; Smithsonian Magazine (1.º de mayo de 2026): \url{https://www.smithsonianmag.com/smart-news/attributed-to-banksy-a-new-statue-of-a-suited-man-180988662/}; MyArtBroker (2026): \url{https://www.myartbroker.com/artist-banksy/articles/banksy-new-sculpture-london-waterloo-place-2026} \\
\end{tabular}

\bigskip

% ---- Introducción contextual ----
\section*{Introducción: la intervención y su contexto}

En la madrugada del 29 de abril de 2026, una escultura de escala monumental apareció sin previo aviso en Waterloo Place, una de las plazas más cargadas de simbolismo imperial del centro de Londres. La obra, instalada secretamente y firmada en su base, representaba a un hombre de traje portando una bandera que, agitada por el viento, le cubría completamente el rostro. El personaje avanza con paso decidido, pero está a punto de caer del borde del plinto sobre el que se sostiene. Al día siguiente, el artista conocido como Banksy reclamó la autoría a través de su cuenta oficial de Instagram.

La ubicación no era arbitraria. Waterloo Place alberga estatuas de figuras emblemáticas del poder imperial y militar británico, entre ellas la del rey Eduardo VII y el Memorial de la Guerra de Crimea. Al insertar su escultura en ese conjunto, Banksy no produjo simplemente una obra de arte público, sino un gesto de intervención crítica en el corazón del paisaje conmemorativo del poder estatal.

% ---- Pregunta 1 ----
\section{¿Qué tipo de poder se representa o cuestiona?}

La escultura pone en cuestión, en primer lugar, el \textbf{poder político-simbólico}: aquel que se ejerce no tanto por la fuerza directa, sino a través de la producción de significados compartidos. La bandera —icono por excelencia de la identidad nacional y la soberanía estatal— es el dispositivo central de la obra. Al convertirla en obstáculo para la visión del propio portador, Banksy interroga el mecanismo por el cual los símbolos nacionales movilizan adhesión emocional y lealtad colectiva.

Hay también una crítica al \textbf{poder institucional y burocrático}. El personaje viste un traje, no un uniforme militar ni un atuendo popular: es el funcionario, el diplomático, el administrador. Como señala el análisis publicado en \textit{Parametric Architecture} (2026), la figura puede leerse como una \enquote{alegoría del sujeto político occidental posicionado entre el discurso de seguridad, el lenguaje diplomático, la idea de interés nacional y un sentido de superioridad histórica}. La obra no apunta exclusivamente al ciudadano ordinario que enarbola una bandera; apunta a quienes toman decisiones desde posiciones de autoridad civil mientras declaran obrar en nombre de principios nacionales que ellos mismos no logran ver con claridad.

Finalmente, la obra dialoga con el \textbf{poder de la memoria histórica}. Al situarse entre monumentos del siglo XIX y comienzos del XX dedicados a guerras y a la expansión imperial, la escultura instaura una conversación anacrónica que relativiza la grandeza conmemorativa que le rodea.

% ---- Pregunta 2 ----
\section{Elementos visuales que construyen la representación}

\subsection*{La figura y su postura}

El hombre avanza con el pie derecho extendido, cruzando ya el límite del plinto. La dirección hacia adelante —gesto típico de las estatuas ecuestes o de los monumentos a héroes— comunica determinación y liderazgo, pero esta convención es inmediatamente desestabilizada por la posición del pie: el sujeto está a punto de caer al vacío. La combinación de movimiento y peligro inminente produce una tensión dramática que es clave para la lectura crítica de la obra.

\subsection*{La bandera}

La bandera no porta ningún símbolo identificable: no hay escudo, no hay colores nacionales específicos. Esta generalidad es deliberada. Como observó un usuario de Instagram citado por CNN (2026), \enquote{la bandera no lleva identidad —ningún país, ninguna lealtad—, solo una forma, lo que hace a la figura universal}. La ausencia de emblema convierte el objeto en un signo puro del concepto de bandera, desnacionalizado y, por ello, potencialmente universal. Su movimiento —impulsada por el viento hacia el rostro del portador— activa la metáfora central: lo que debería ser estandarte de visión colectiva se convierte en venda.

\subsection*{El traje}

La vestimenta de civil —saco, corbata, pantalón de vestir— ancla al personaje en el ámbito del poder administrativo y de representación formal, no en el de la guerra ni en el del activismo popular. Esa elección indica que la crítica no recae sobre el soldado razo ni sobre el ciudadano entusiasta, sino sobre las élites que administran los símbolos nacionales.

\subsection*{El plinto y el entorno}

El plinto imita el lenguaje formal de los monumentos históricos circundantes. Banksy adopta la misma gramática visual —base elevada, materiales pétreos, figura en escala monumental— para subvertirla desde adentro. Donde cada monumento de Waterloo Place proyecta permanencia y compostura, el suyo anuncia inestabilidad e inminente colapso. El crítico Alex Dale, citado por Wikipedia (2026), describió la pieza como dotada de \enquote{estética de ornamento de jardín y complejidad moral de un acertijo}.

% ---- Pregunta 3 ----
\section{Mensaje e interpretación: ¿qué hace la imagen con el poder?}

La escultura no legitima el poder; tampoco lo denuncia desde una perspectiva de víctima directa. Su operación es la de la \textbf{ironía y la subversión}. Banksy toma las formas del monumento público —el plinto, el traje formal, el gesto de avanzar— y las vacía de su función legitimadora al introducir un solo elemento perturbador: la ceguera provocada por el propio símbolo que el sujeto abraza.

El mensaje puede sintetizarse en lo que la obra de arte llama ``patriotismo ciego'' (\textit{blind patriotism}), pero su alcance es más amplio. Sugiere que la lealtad sin crítica a los símbolos colectivos —bandera, nación, estado— no moviliza genuinamente hacia ningún destino: coloca a quien la practica al borde del precipicio. El movimiento hacia adelante, que en los monumentos convencionales significa progreso, aquí significa riesgo de caída. La confianza ciega en el símbolo nacional reemplaza a la deliberación y a la visión política real.

El \textit{Blind Patriotism} no propone una alternativa explícita ni articula un programa político. Opera en el registro de la pregunta más que de la respuesta: ¿hacia dónde lleva marchar bajo una bandera que te impide ver?

% ---- Pregunta 4 ----
\section{Relación con las ideas de Peter Burke en ``Poder y Protesta''}

En el capítulo cuarto de \textit{Visto y no visto. El uso de la imagen como documento histórico} (2001), Peter Burke examina el papel de las imágenes tanto como instrumentos de afirmación del poder como vehículos de protesta y subversión política. Una de sus tesis centrales es que \enquote{las imágenes son una forma importante de documento histórico} que va más allá del mero registro: \enquote{reflejan un testimonio ocular} de las mentalidades, los valores y los conflictos de su época. Burke observa que, desde la Antigüedad, los gobernantes han utilizado la imagen para legitimar su autoridad —mediante convenciones visuales como la vestimenta, la pose, la escala, los gestos—, y que los movimientos de protesta han respondido recurriendo a las mismas imágenes, pero invirtiéndolas o resignificándolas.

Esta dinámica es exactamente la que Banksy pone en práctica. La escultura de Waterloo Place no inventa un lenguaje propio: toma prestado el lenguaje del monumento público —ese dispositivo visual por excelencia de la autorepresentación del poder estatal— y lo desvía. Burke señala que en las imágenes de poder conviven elementos que \enquote{buscan legitimar, ironizar, subvertir o denunciar}, y que el historiador debe distinguir entre la intención del productor, la lectura del espectador coetáneo y la interpretación posterior. En el caso de Banksy, los tres niveles coinciden de manera poco frecuente: la intención subversiva es explícita, la recepción inmediata fue masivamente crítica hacia el poder, y la lectura histórica futura probablemente la inscriba en el debate sobre el auge del nacionalismo populista de los años 2020.

Burke también advierte que las imágenes no son testimonios neutros: el artista \enquote{selecciona qué aspectos del mundo real va a plasmar} y no renuncia a una puesta en escena, a convenciones culturales ni a un posicionamiento moral o político. La escultura de Banksy cumple exactamente este criterio: su aparente objetividad formal —la figura realista, el plinto sobrio— encubre una fortísima toma de partido que solo se revela en el detalle de la bandera que ciega.

% ---- Pregunta 5 ----
\section{Límites y riesgos de usar esta imagen como documento histórico}

\subsection*{La intención autoral como trampa hermenéutica}

Banksy construye deliberadamente un mensaje inequívoco: la ironía está codificada de manera casi didáctica. Eso es, al mismo tiempo, su fortaleza comunicativa y su límite historiográfico. Como advierte Burke, cuanto más explícita es la intención propagandística o subversiva de una imagen, más opaca resulta para el análisis de las condiciones sociales que la rodean. La escultura nos dice más sobre el posicionamiento crítico de su autor que sobre el estado real del patriotismo o del nacionalismo en la sociedad británica de 2026.

\subsection*{La autoría anónima y la mediación institucional}

La identidad no confirmada de Banksy introduce una incertidumbre documental relevante. La firma en la base y la publicación en Instagram son pruebas de reivindicación, no de autoría en sentido estricto. Más significativo aún: la reacción de las instituciones —el alcalde de Londres y el Westminster City Council decidieron no retirar la pieza— forma parte del documento. La escultura fue absorbida rápidamente por el sistema que pretendía criticar: el poder estatal la protegió, la preservó y la convirtió en atractivo turístico. Esta cooptación institucional limita su potencial subversivo real y complica su lectura como gesto de protesta genuina frente al poder.

\subsection*{La circulación digital como distorsión}

Las imágenes de la escultura se difundieron globalmente a través de redes sociales antes de que ningún medio de comunicación pudiera contextualizar la obra. Esa circulación viralizó una interpretación determinada —``el hombre cegado por la bandera como metáfora del patriotismo irracional''— antes de que pudieran emerger lecturas alternativas. Para Burke, el historiador que usa imágenes debe preguntarse siempre quién produce la imagen, para quién, en qué contexto y con qué mediaciones. En este caso, la mediación digital masiva homogeneizó la interpretación y dificultó lecturas disidentes o matizadas.

\subsection*{El riesgo de la iconización}

Existe, finalmente, el riesgo de que la escultura sea leída en el futuro como un documento del ``espíritu de la época'' cuando en realidad es una intervención polémica de un autor con posicionamiento político definido. Toda imagen que alcanza notoriedad corre el riesgo de ser naturalizada como evidencia histórica de aquello que solo es su perspectiva sobre la historia.

% ---- Conclusión ----
\section*{Conclusión}

La escultura instalada por Banksy en Waterloo Place el 29 de abril de 2026 es un caso ejemplar de lo que Peter Burke llama el uso de la imagen como intervención en los debates sobre el poder. La obra no solo representa el poder: lo interpela, lo ironiza y lo desestabiliza desde adentro, apropiándose de sus propios lenguajes formales. Situada en uno de los escenarios más cargados de simbolismo imperial de Londres, y en un momento de intensificación de los nacionalismos populistas a escala global, la escultura actúa como un contra-monumento que hace visible la lógica ciega de la adhesión simbólica.

Al mismo tiempo, su análisis como documento histórico exige la cautela metodológica que Burke prescribe: toda imagen está construida desde una perspectiva, toda imagen puede ser leída de múltiples maneras, y toda imagen puede ser apropiada —incluso por aquellos a quienes critica— para fines distintos de los que su autor concibió.

% ---- Bibliografía ----
\section*{Bibliografía y fuentes}

\begin{enumerate}
    \item Burke, Peter. \textit{Visto y no visto. El uso de la imagen como documento histórico}. Traducción de Teófilo de Lozoya. Barcelona: Crítica, 2001. [Cap. 4: ``Poder y Protesta'']

    \item NPR. \enquote{Banksy confirms new statue installed in central London is his work.} 2 de mayo de 2026. \url{https://www.npr.org/2026/05/02/nx-s1-5808871/banksy-statue-london}

    \item CNN. \enquote{Banksy's new flag-wielding London statue satirizes blind patriotism.} 1.º de mayo de 2026. \url{https://www.cnn.com/2026/05/01/style/banksy-new-statue-london-intl-hnk}

    \item Smithsonian Magazine. \enquote{Attributed to Banksy, a New Statue of a Suited Man, Blinded by a Flag.} 1.º de mayo de 2026. \url{https://www.smithsonianmag.com/smart-news/attributed-to-banksy-a-new-statue-of-a-suited-man-180988662/}

    \item Wikipedia. \enquote{Banksy statue.} Consultada el 8 de mayo de 2026. \url{https://en.wikipedia.org/wiki/Banksy_statue}

    \item MyArtBroker. \enquote{Banksy's New Sculpture in London: A Suited Figure, a Flag, and a Fall.} 2026. \url{https://www.myartbroker.com/artist-banksy/articles/banksy-new-sculpture-london-waterloo-place-2026}

    \item Parametric Architecture. \enquote{Flag and Void: A Historical and Contemporary Reading of Banksy's Sculpture at Waterloo Place.} Mayo de 2026. \url{https://parametric-architecture.com/banksys-sculpture-at-waterloo-place/}

    \item Class Art Biarritz. \enquote{Blind Patriotism, a new sculpture by Banksy in London.} 2026. \url{https://www.classartbiarritz.net/post/blind-patriotism-a-new-sculpture-by-banksy-in-london}

    \item CBC News. \enquote{Street artist Banksy takes credit for London statue.} 30 de abril de 2026. \url{https://www.cbc.ca/news/entertainment/banksy-statue-london-9.7183201}
\end{enumerate}

\end{document}